\subsection*{Случай наличия обновлений на подотрезке}

\textbf{Идея отложенных операций} 
Применим следующий трюк: если нам нужно обновить значение подотрезка, то мы отложим эту операцию до востребования. Запомним в соответствующих вершине/блоке, что мы должны в будущем сделать обновление в потомков. Если нам вдруг понадобится значение, лежащее в вершине/блоке, то мы посчитаем его так, как будто мы обновили отрезок, а напоминание о том, что надо сделать со всеми элементами, мы "протолкнем"в потомков.

\Def Операцией группового обновления ch будем называть запрос, в
ходе которого применяется на подотрезке операция $ch(\cdot,\,val)$.

Свойства на операцию ch, если op — операция запроса на отрезке.

\begin{enumerate}
    \item Существование нейтрального элемента
    \item Ассоциативность: $ch(a,\,ch(b,\,c))\,=\,ch(ch(a,\,b),\,c)$
    \item Дистрибутивность: $ch(op(a,\,b),\,c)\,=\,op(ch(a,\,c),\,ch(b,\,c))$
\end{enumerate}

\Def Несогласованность -- величина, с которой надо выполнить операцию op, чтобы получить корректный результат.

В каждый момент времени в узле необязательно лежит истинное значение op на отрезке. Но к моменту запроса будем осуществлять проталкивание несогласованности.

\begin{lstlisting}
void push(int node) {
  tree[2 * node + 1].lazy = ch(tree[2 * node + 1].lazy, tree[node].lazy);
  tree[2 * node + 2].lazy = ch(tree[2 * node + 2].lazy, tree[node].lazy);

  tree[node].d = ch_neutral
}
\end{lstlisting}

Вызывать push надо каждый раз, когда мы обращаемся к вершине, которая могла
быть изменена

\subsubsection{Запрос на группового обновления}

\begin{lstlisting}
void update(int node, int a, int b, T val) {
  l = tree[node].left
  r = tree[node].right

  if ([l, r) is subset [a, b)) {
    tree[node].lazy = ch(tree[node].lazy, val);
    return;
  }
  push(node);

  update(2 * node + 1, a, b, val);
  update(2 * node + 2, a, b, val);
  tree[node].val = op(
    ch(tree[2 * node + 1].val, tree[2 * node + 1].lazy),
    ch(tree[2 * node + 2].val, tree[2 * node + 2].lazy),
  );
}
\end{lstlisting}

\Alg
\begin{enumerate}
    \item Первая часть посвящена разбиению на объединение подотрезков. Если подмножество, то накапливаем несогласованность.
    \item Проталкивание.
    \item Обновляем детей.
    \item Обновляем результат. Заметим, что у нас несогласованность уже нейтральная, так что надо просто собрать op с «истинных» значений у детей.
\end{enumerate}

\subsection{Ответ на запрос}

Выполняется аналогично групповому обновлению

\begin{lstlisting}
int query(int node, int a, int b) {
  l = tree[node].left
  r = tree[node].right

  if ([l, r) is subset [a, b)) {
    return ch(tree[node].val, tree[node].lazy);
  }
  push(node);

  answer = op(
    query(2 * node + 1, a, b),
    query(2 * node + 2, a, b),
  );

  tree[node].val = op(
    ch(tree[2 * node + 1].val, tree[2 * node + 1].lazy),
    ch(tree[2 * node + 2].val, tree[2 * node + 2].lazy),
  );
  return answer;
}
\end{lstlisting}

\pagebreak